%
\section{Experiencias musicales digitales e interactivas}
\subsection{Transformación del consumo musical digital}
Durante la última década, el consumo musical ha experimentado una transformación significativa debido a la expansión de las plataformas digitales. A diferencia del modelo tradicional basado en la compra y propiedad de contenidos individuales, los servicios de \emph{streaming} permiten acceder, generalmente mediante una suscripción, a amplios catálogos musicales sin pagar por cada reproducción. Este cambio ha incrementado la cantidad y diversidad de música consumida y ha modificado las formas de descubrimiento musical, ya que los usuarios recurren a servicios de \emph{streaming}, plataformas audiovisuales y redes sociales para encontrar nuevas canciones, artistas y videoclips \cite{Datta2018}.
\subsection{El videoclip como experiencia audiovisual}
El videoclip es un formato audiovisual que integra música, imágenes y recursos narrativos para comunicar y ampliar el significado de una canción. La relación entre la letra, las escenas y la propuesta visual puede influir en la interpretación y en la respuesta emocional del público \cite{saiyidatul2024narrative}. Con su incorporación a plataformas digitales, el formato también ha evolucionado hacia videoclips interactivos, en los cuales la participación del usuario interviene en la presentación y construcción de la narrativa audiovisual \cite{barboza2015digital}.
\subsection{Experiencias musicales en realidad virtual}
Las experiencias musicales en realidad virtual combinan música, entornos tridimensionales y posibilidades de exploración que permiten al usuario percibir el evento desde perspectivas difíciles y distintas a las de un concierto físico. La posibilidad de interactuar con escenarios y recursos visuales imposibles de reproducir en la realidad constituye uno de sus principales atractivos \cite{onderdijk2023concert}. Además, el uso de visores y la navegación activa entre diferentes puntos de vista pueden incrementar la presencia, el estado de flujo, la satisfacción y la intención de volver a participar en conciertos virtuales \cite{lo2023immersive}.

La literatura sobre conciertos virtuales también aborda recursos específicos para enriquecer la experiencia de audiencia. Entre estos se encuentran el diseño de múltiples vistas para escoger perspectivas del escenario \cite{kim2022multiview}, el análisis de formatos de visualización y géneros musicales en conciertos con avatares \cite{lee2023viewing}, las reacciones colectivas artificiales basadas en conversación \cite{lee2025concert} y las interacciones sociales para presentaciones en vivo dentro de entornos virtuales \cite{yan2020social}. También existen antecedentes de producciones sociales de conciertos en VR con audiencias masivas, como \emph{Welcome to the Other Side} \cite{cacciuttolo2021welcome}. En el contexto latinoamericano, el diseño gráfico de experiencias de realidad virtual aplicadas a la presentación de música urbana muestra la relación entre identidad visual, escena musical y entorno inmersivo \cite{rivas2025diseno}.
\subsection{Videojuegos musicales e interacción con la música}
Los videojuegos musicales convierten la música en un componente central de la interacción, en lugar de limitarla únicamente a un acompañamiento. En ellos, las acciones del jugador pueden relacionarse con el ritmo, la interpretación y la producción de respuestas sonoras o visuales. Elementos propios del juego, como desafíos, puntuaciones y movimientos vinculados con la respuesta audiovisual, pueden fomentar una escucha activa al requerir que el usuario preste atención a características como el ritmo, la métrica o la estructura musical \cite{mandanici2018good}. De esta misma forma, entre los principios de interacción identificados en los videojuegos musicales se encuentran las partituras activas, las acciones rítmicas, la cuantización, la sinestesia, la interpretación como forma de juego y la manipulación libre del sonido \cite{pichlmair2007levels}.
\subsection{Generación procedural de contenido}
La generación procedural de contenido permite producir elementos de una experiencia mediante reglas, algoritmos o sistemas de selección, lo cual puede apoyar la variación y adaptación del contenido interactivo. En videojuegos, este enfoque puede orientarse por la experiencia del usuario cuando el sistema utiliza indicadores afectivos o de comportamiento para guiar la generación \cite{yannakakis2011experience}. Sin embargo, la generación procedural no sustituye necesariamente la intención autoral: también puede entenderse mediante metáforas de diseño que ayudan a decidir qué aspectos controla el sistema y cuáles permanecen bajo control del diseñador \cite{khaled2013design}. Desde esta perspectiva, una aproximación autoral permite usar contenido procedural sin perder la dirección expresiva de la experiencia \cite{craveirinha2015designing}.

% 
\section{Desarrollo de experiencias de realidad virtual en Unity}
\subsection{Realidad virtual y entornos virtuales interactivos}
La realidad virtual permite situar al usuario dentro de un entorno generado digitalmente mediante tecnologías que capturan sus movimientos, presentan estímulos a sus sentidos y actualizan la experiencia en tiempo real. De esta forma, un entorno virtual interactivo no se limita a representar un espacio tridimensional, sino que responde a las acciones del usuario y permite establecer una relación continua entre sus movimientos, las entradas del sistema y los cambios producidos dentro del entorno \cite{merienne2016virtual}. La interacción, junto con la inmersión y la presencia, constituye así un componente relevante de la experiencia en realidad virtual \cite{merienne2016virtual}.
\subsection{Unity}
Unity es un motor de desarrollo en tiempo real utilizado para crear juegos y aplicaciones en dos y tres dimensiones para diferentes plataformas \cite{unity2026engine}. Su entorno organiza el contenido por medio de escenas compuestas por objetos, componentes y recursos, mientras que su ecosistema permite incorporar sistemas de física, renderizado, audio, programación e interacción. Para el desarrollo de realidad virtual, Unity dispone además de herramientas y paquetes especializados que permiten integrar dispositivos XR y construir sistemas de interacción dentro de los entornos virtuales \cite{unity2026engine,unity2026vroverview}.
\subsection{C\#}
C\# es un lenguaje de programación de uso general, orientado a objetos y con tipado fuerte, utilizado por Unity para desarrollar la lógica de sus aplicaciones \cite{microsoft2026csharp,unity2026scripting}. Mediante \emph{scripts} escritos en C\# es posible controlar el comportamiento de los objetos de una escena, responder a las entradas del usuario, ejecutar eventos y modificar elementos de la aplicación durante su ejecución. En Unity, estos scripts pueden convertirse en componentes por medio de clases derivadas de \emph{MonoBehaviour}, para relacionar el comportamiento programado con los \emph{GameObjects} \cite{unity2026scripting}.
\subsection{Plantilla de realidad virtual y escena de muestra de Unity}
La plantilla de realidad virtual de Unity proporciona un punto de partida preconfigurado para desarrollar aplicaciones VR, incluyendo ajustes del proyecto, paquetes necesarios y una escena de muestra con componentes y recursos preparados para demostrar interacciones habituales \cite{unity2026vroverview}. Esta escena permite observar ejemplos funcionales de interacción, como manipulación de objetos y otras mecánicas propias de realidad virtual, por lo que puede utilizarse como referencia para comprender la configuración de una experiencia antes de desarrollar sistemas propios \cite{unity2026vrtemplate}.
\subsection{\emph{XR Origin} e interacción mediante controladores de movimiento}
El \emph{XR Origin} representa el sistema de referencia del usuario dentro de una escena de Realidad Extendida (XR por sus siglas en inglés) y agrupa elementos asociados con el seguimiento de la cabeza, la cámara y los dispositivos de entrada. A partir de la información de posicionamiento y rotación proporcionada por el visor y los controladores, Unity actualiza su representación dentro del entorno virtual \cite{unity2026xritoolkit}. El \emph{XR Interaction Toolkit} complementa este sistema mediante componentes Interactor e Interactable, que permiten implementar acciones como seleccionar, tomar y manipular objetos, además de interactuar con interfaces y sistemas de locomoción utilizando las entradas de los controladores de movimiento \cite{unity2026xritoolkit,unity2026xrinput}.

% 
\section{Presencia y evaluación de la experiencia}
\subsection{Presencia en realidad virtual}
La presencia se relaciona con la sensación subjetiva de encontrarse dentro del entorno virtual y responder a este como si fuera un espacio experimentado directamente. En aplicaciones de realidad virtual, esta percepción depende de factores técnicos, perceptuales e interactivos, como la calidad de los estímulos, la correspondencia entre movimiento físico y respuesta visual, el grado de involucramiento y la coherencia del entorno \cite{merienne2016virtual}. En experiencias musicales inmersivas, la presencia es relevante porque influye en cómo el usuario interpreta el espacio, la cercanía con la presentación y el nivel de participación dentro de la experiencia \cite{onderdijk2023concert,lo2023immersive}.
\subsection{\emph{I-Group Presence Questionnaire}}
El \emph{I-Group Presence Questionnaire} (IPQ) es un instrumento utilizado para evaluar la presencia en entornos virtuales mediante dimensiones como presencia espacial, involucramiento y realismo experimentado. La versión revisada en inglés del IPQ analiza sus propiedades psicométricas y proporciona una base para su uso en estudios de realidad virtual \cite{berkman2021igroup}. Asimismo, la adaptación al español de los cuestionarios ITQ e IPQ permite considerar su aplicación en población hispanohablante y relacionar presencia, tendencia inmersiva, sugestionabilidad y malestar por realidad virtual \cite{castro2026immersive}.
\subsection{Evaluación de experiencias musicales en VR}
La evaluación de experiencias musicales en VR puede considerar tanto variables propias de la presencia como indicadores específicos de la experiencia musical, tales como satisfacción, flujo, intención de volver a participar, interacción social y percepción del formato de visualización \cite{lo2023immersive,lee2023viewing}. En este tipo de experiencias, las decisiones de navegación, perspectiva, avatarización e interacción colectiva pueden modificar la forma en que el usuario vive el evento y, por tanto, deben vincularse con instrumentos de evaluación que midan su experiencia subjetiva \cite{lee2025concert,kim2022multiview,yan2020social}.

\section{Arquitectura dirigida por eventos y secuenciación inmersiva}
\subsection{Eventos y coordinación de sistemas interactivos}
Una arquitectura dirigida por eventos organiza el comportamiento de un sistema a partir de eventos que notifican cambios, acciones o estados relevantes. Este enfoque permite desacoplar los componentes que producen eventos de los componentes que reaccionan a ellos, lo que resulta útil en experiencias interactivas donde audio, visuales, animaciones e interacción del usuario deben coordinarse durante la ejecución. En entornos virtuales, la visualización inmersiva de flujos de eventos y datos evidencia cómo los eventos pueden representarse, organizarse y analizarse dentro de VR \cite{oberhauser2023vr}.
\subsection{Eventos en entornos virtuales y mixtos}
Los sistemas de realidad virtual y aumentada requieren coordinar estímulos, entradas y respuestas en tiempo real. En trabajos sobre entornos mixtos y aumentados se han explorado técnicas para vincular acciones, percepción e interacción multisensorial mediante sistemas que reaccionan a eventos del usuario y del entorno \cite{edmunds2006mixed}. Esta lógica es pertinente para un videoclip musical en VR, ya que la experiencia debe sincronizar cambios visuales, objetos interactivos y momentos narrativos con la progresión temporal de la canción.
\subsection{Secuenciación autoral para videoclips musicales en VR}
La secuenciación de un videoclip musical en VR puede combinar eventos temporizados, recursos interactivos y decisiones autorales. A diferencia de una experiencia generada completamente de forma procedural, el sistema puede utilizar eventos para organizar cambios predefinidos que preservan la intención narrativa y estética de la pieza musical \cite{khaled2013design,craveirinha2015designing}. Este enfoque permite integrar principios de interacción musical y diseño inmersivo sin desplazar el control creativo necesario para mantener coherencia entre música, escena y experiencia del usuario \cite{pichlmair2007levels,saiyidatul2024narrative}.

% 
\section{Presencia y evaluación de la experiencia}
\subsection{}
\subsection{}
\subsection{}
\subsection{}


% 
\section{}
\subsection{}
\subsection{}
\subsection{}
\subsection{}